% ============================================================================
%  Findings section for short.tex
%  This file is a TeX fragment, intended to be included with:
%  % ============================================================================
%  Findings section for short.tex
%  This file is a TeX fragment, intended to be included with:
%  % ============================================================================
%  Findings section for short.tex
%  This file is a TeX fragment, intended to be included with:
%  % ============================================================================
%  Findings section for short.tex
%  This file is a TeX fragment, intended to be included with:
%  \input{findings_global_formation}
% ============================================================================

\section{Findings: Global AI Startup Formation}
\label{sec:findings_global_formation}

The main descriptive pattern is simple: AI has become a much larger share of
new startup formation. The harder part is interpretation. Recent founding
cohorts are still being indexed, so raw counts understate recent formation.
For that reason, the formation result is stated primarily as a \emph{share of
observed firms founded in the same year}, not as the raw number of AI firms.

\subsection{AI's share of new firms rises sharply}

In the commercial population, where founding years are nearly complete, the AI
share rises from $3.0\%$ in 2000 to $10.6\%$ in 2015, $18.9\%$ in 2021, and
$59.0\%$ in 2025 (Table~\ref{tab:formation_selected_years}). The largest break
in the series is between the 2022 and 2023 founding cohorts, when the AI share
jumps from $25.5\%$ to $43.5\%$. That timing is consistent with the public
release of consumer generative AI in late 2022, but we do not interpret the
jump causally. The same cohorts are also affected by indexing lag, and the
paper does not have a design that separates the technology shock from the
change in observability.

\begin{table}[htbp]
\centering\small
\caption{AI share of observed startup formation, selected founding cohorts.
The sample is the commercial population, where founding years are nearly
complete.}
\label{tab:formation_selected_years}
\begin{tabular}{@{}rrrr@{}}
\toprule
\textbf{Founded year} & \textbf{Observed firms} & \textbf{AI firms} &
\textbf{AI share} \\
\midrule
2000 & $28{,}963$ & $874$ & $3.0\%$ \\
2010 & $41{,}513$ & $1{,}996$ & $4.8\%$ \\
2015 & $54{,}150$ & $5{,}764$ & $10.6\%$ \\
2018 & $50{,}542$ & $8{,}038$ & $15.9\%$ \\
2021 & $35{,}978$ & $6{,}812$ & $18.9\%$ \\
2022 & $22{,}800$ & $5{,}818$ & $25.5\%$ \\
2023 & $17{,}947$ & $7{,}812$ & $43.5\%$ \\
2024 & $10{,}562$ & $5{,}534$ & $52.4\%$ \\
2025 & $2{,}531$ & $1{,}492$ & $59.0\%$ \\
\bottomrule
\end{tabular}
\end{table}

The count series tells a different story. The observed number of AI firms peaks
in 2018 at $8{,}038$ and falls after that. We do not read this as a decline in
AI entrepreneurship. The total observed denominator falls from $54{,}150$ firms
in the 2015 cohort to $2{,}531$ in 2025, because recent firms have not yet had
time to enter commercial registers. The useful fact is therefore not that
observed AI counts fall after 2018, but that AI's share continues to rise even
as the recent denominator becomes incomplete.

\subsection{The trend is not driven by one source}

The rise appears in both commercial databases. Database A rises from $2.7\%$
AI in the 2000 cohort to $63.7\%$ in 2025. Database B rises from $3.6\%$ to
$58.3\%$ over the same window. The levels differ somewhat, but the direction
does not. This matters because a trend visible in only one database could
reflect a vendor-specific tagging or coverage decision. Here, both sources
show the same broad movement.

\subsection{The trend survives a stricter selection rule}

Recent cohorts are not only smaller; they are also more selected. A young firm
is especially likely to appear in a commercial database if it has already done
something visible, such as raising financing. If AI firms are more likely to be
visible early, the rising AI share could partly reflect selection into the
database rather than startup formation itself.

We check this by holding the selection rule fixed: among firms with recorded
financing, the AI share still rises from $5.9\%$ in 2005 to $15.6\%$ in 2015,
$22.0\%$ in 2021, and $68.3\%$ in 2025. The level is higher because financed
firms are more AI-intensive, but the time pattern remains. This does not solve
all selection problems, but it rules out the simplest concern that the trend is
only an artifact of looser recent coverage.

\subsection{Commercial databases miss AI-heavy firms}

The formation series above uses the commercial population because it has the
founding-year metadata needed for a time trend. The broader dataset also
reveals what those commercial sources miss. Firms found outside both
commercial databases are $23.8\%$ AI-related, compared with $12.4\%$ in
database A and $10.7\%$ in database B. The unlisted layer is therefore roughly
twice as AI-intensive as the conventional evidence base.

We do not use the unlisted layer to extend the time series, because its
founding-year coverage is thin. Its value is different: it shows the direction
of coverage bias. If the firms missing from commercial databases are more
likely to be AI-related, then commercial-only studies will understate AI entry.
Section~\ref{sec:checks} tests whether that gap survives alternative AI
definitions, source exclusions, independent labels, and classifier-error
correction.

\subsection{The global pattern is broad growth with cautious geography}

AI intensity rises across every major region we can measure. Between the
2010--2016 and 2020--2025 cohorts, North America's AI share rises from $9.1\%$
to $29.2\%$, Europe's from $7.3\%$ to $21.4\%$, East Asia's from $12.3\%$ to
$25.1\%$, and South and Southeast Asia's from $7.2\%$ to $18.2\%$. Among large
country samples, recent AI intensity is highest in Israel ($37.2\%$), the
United States ($29.4\%$), China ($28.5\%$), Canada ($27.3\%$), and South Korea
($26.1\%$).

The concentration series is more delicate. Through 2017, AI entrepreneurship
appears to become more geographically dispersed: the number of observed
countries with AI firms rises from $59$ in 2000 to $103$ in 2017, and the U.S.\
share of observed AI firms falls from $41.3\%$ to $36.8\%$. After 2020 the
series reverses sharply, with the U.S.\ reaching $85.0\%$ of observed AI firms
in the 2025 cohort. We treat that reversal as a coverage warning rather than a
substantive geography result, because the newest non-U.S.\ firms are indexed
more slowly.

\subsection{AI growth is broad across sectors}

The sector pattern is also one of broad intensification rather than a single
vertical driving the result. Data and analytics rises from $35.9\%$ AI in the
2010--2016 cohort to $71.4\%$ in 2020--2025. Biotech and pharma rises from
$23.2\%$ to $60.0\%$, software and IT from $9.5\%$ to $29.3\%$, and healthcare
from $8.9\%$ to $26.6\%$. Lower-intensity sectors also rise: food and agtech,
consumer and lifestyle, and e-commerce and retail all more than double their
AI shares. The ordering changes less than the levels. Sectors that were already
AI-intensive remain near the top, while nearly every sector moves upward.

% ============================================================================

\section{Findings: Global AI Startup Formation}
\label{sec:findings_global_formation}

The main descriptive pattern is simple: AI has become a much larger share of
new startup formation. The harder part is interpretation. Recent founding
cohorts are still being indexed, so raw counts understate recent formation.
For that reason, the formation result is stated primarily as a \emph{share of
observed firms founded in the same year}, not as the raw number of AI firms.

\subsection{AI's share of new firms rises sharply}

In the commercial population, where founding years are nearly complete, the AI
share rises from $3.0\%$ in 2000 to $10.6\%$ in 2015, $18.9\%$ in 2021, and
$59.0\%$ in 2025 (Table~\ref{tab:formation_selected_years}). The largest break
in the series is between the 2022 and 2023 founding cohorts, when the AI share
jumps from $25.5\%$ to $43.5\%$. That timing is consistent with the public
release of consumer generative AI in late 2022, but we do not interpret the
jump causally. The same cohorts are also affected by indexing lag, and the
paper does not have a design that separates the technology shock from the
change in observability.

\begin{table}[htbp]
\centering\small
\caption{AI share of observed startup formation, selected founding cohorts.
The sample is the commercial population, where founding years are nearly
complete.}
\label{tab:formation_selected_years}
\begin{tabular}{@{}rrrr@{}}
\toprule
\textbf{Founded year} & \textbf{Observed firms} & \textbf{AI firms} &
\textbf{AI share} \\
\midrule
2000 & $28{,}963$ & $874$ & $3.0\%$ \\
2010 & $41{,}513$ & $1{,}996$ & $4.8\%$ \\
2015 & $54{,}150$ & $5{,}764$ & $10.6\%$ \\
2018 & $50{,}542$ & $8{,}038$ & $15.9\%$ \\
2021 & $35{,}978$ & $6{,}812$ & $18.9\%$ \\
2022 & $22{,}800$ & $5{,}818$ & $25.5\%$ \\
2023 & $17{,}947$ & $7{,}812$ & $43.5\%$ \\
2024 & $10{,}562$ & $5{,}534$ & $52.4\%$ \\
2025 & $2{,}531$ & $1{,}492$ & $59.0\%$ \\
\bottomrule
\end{tabular}
\end{table}

The count series tells a different story. The observed number of AI firms peaks
in 2018 at $8{,}038$ and falls after that. We do not read this as a decline in
AI entrepreneurship. The total observed denominator falls from $54{,}150$ firms
in the 2015 cohort to $2{,}531$ in 2025, because recent firms have not yet had
time to enter commercial registers. The useful fact is therefore not that
observed AI counts fall after 2018, but that AI's share continues to rise even
as the recent denominator becomes incomplete.

\subsection{The trend is not driven by one source}

The rise appears in both commercial databases. Database A rises from $2.7\%$
AI in the 2000 cohort to $63.7\%$ in 2025. Database B rises from $3.6\%$ to
$58.3\%$ over the same window. The levels differ somewhat, but the direction
does not. This matters because a trend visible in only one database could
reflect a vendor-specific tagging or coverage decision. Here, both sources
show the same broad movement.

\subsection{The trend survives a stricter selection rule}

Recent cohorts are not only smaller; they are also more selected. A young firm
is especially likely to appear in a commercial database if it has already done
something visible, such as raising financing. If AI firms are more likely to be
visible early, the rising AI share could partly reflect selection into the
database rather than startup formation itself.

We check this by holding the selection rule fixed: among firms with recorded
financing, the AI share still rises from $5.9\%$ in 2005 to $15.6\%$ in 2015,
$22.0\%$ in 2021, and $68.3\%$ in 2025. The level is higher because financed
firms are more AI-intensive, but the time pattern remains. This does not solve
all selection problems, but it rules out the simplest concern that the trend is
only an artifact of looser recent coverage.

\subsection{Commercial databases miss AI-heavy firms}

The formation series above uses the commercial population because it has the
founding-year metadata needed for a time trend. The broader dataset also
reveals what those commercial sources miss. Firms found outside both
commercial databases are $23.8\%$ AI-related, compared with $12.4\%$ in
database A and $10.7\%$ in database B. The unlisted layer is therefore roughly
twice as AI-intensive as the conventional evidence base.

We do not use the unlisted layer to extend the time series, because its
founding-year coverage is thin. Its value is different: it shows the direction
of coverage bias. If the firms missing from commercial databases are more
likely to be AI-related, then commercial-only studies will understate AI entry.
Section~\ref{sec:checks} tests whether that gap survives alternative AI
definitions, source exclusions, independent labels, and classifier-error
correction.

\subsection{The global pattern is broad growth with cautious geography}

AI intensity rises across every major region we can measure. Between the
2010--2016 and 2020--2025 cohorts, North America's AI share rises from $9.1\%$
to $29.2\%$, Europe's from $7.3\%$ to $21.4\%$, East Asia's from $12.3\%$ to
$25.1\%$, and South and Southeast Asia's from $7.2\%$ to $18.2\%$. Among large
country samples, recent AI intensity is highest in Israel ($37.2\%$), the
United States ($29.4\%$), China ($28.5\%$), Canada ($27.3\%$), and South Korea
($26.1\%$).

The concentration series is more delicate. Through 2017, AI entrepreneurship
appears to become more geographically dispersed: the number of observed
countries with AI firms rises from $59$ in 2000 to $103$ in 2017, and the U.S.\
share of observed AI firms falls from $41.3\%$ to $36.8\%$. After 2020 the
series reverses sharply, with the U.S.\ reaching $85.0\%$ of observed AI firms
in the 2025 cohort. We treat that reversal as a coverage warning rather than a
substantive geography result, because the newest non-U.S.\ firms are indexed
more slowly.

\subsection{AI growth is broad across sectors}

The sector pattern is also one of broad intensification rather than a single
vertical driving the result. Data and analytics rises from $35.9\%$ AI in the
2010--2016 cohort to $71.4\%$ in 2020--2025. Biotech and pharma rises from
$23.2\%$ to $60.0\%$, software and IT from $9.5\%$ to $29.3\%$, and healthcare
from $8.9\%$ to $26.6\%$. Lower-intensity sectors also rise: food and agtech,
consumer and lifestyle, and e-commerce and retail all more than double their
AI shares. The ordering changes less than the levels. Sectors that were already
AI-intensive remain near the top, while nearly every sector moves upward.

% ============================================================================

\section{Findings: Global AI Startup Formation}
\label{sec:findings_global_formation}

The main descriptive pattern is simple: AI has become a much larger share of
new startup formation. The harder part is interpretation. Recent founding
cohorts are still being indexed, so raw counts understate recent formation.
For that reason, the formation result is stated primarily as a \emph{share of
observed firms founded in the same year}, not as the raw number of AI firms.

\subsection{AI's share of new firms rises sharply}

In the commercial population, where founding years are nearly complete, the AI
share rises from $3.0\%$ in 2000 to $10.6\%$ in 2015, $18.9\%$ in 2021, and
$59.0\%$ in 2025 (Table~\ref{tab:formation_selected_years}). The largest break
in the series is between the 2022 and 2023 founding cohorts, when the AI share
jumps from $25.5\%$ to $43.5\%$. That timing is consistent with the public
release of consumer generative AI in late 2022, but we do not interpret the
jump causally. The same cohorts are also affected by indexing lag, and the
paper does not have a design that separates the technology shock from the
change in observability.

\begin{table}[htbp]
\centering\small
\caption{AI share of observed startup formation, selected founding cohorts.
The sample is the commercial population, where founding years are nearly
complete.}
\label{tab:formation_selected_years}
\begin{tabular}{@{}rrrr@{}}
\toprule
\textbf{Founded year} & \textbf{Observed firms} & \textbf{AI firms} &
\textbf{AI share} \\
\midrule
2000 & $28{,}963$ & $874$ & $3.0\%$ \\
2010 & $41{,}513$ & $1{,}996$ & $4.8\%$ \\
2015 & $54{,}150$ & $5{,}764$ & $10.6\%$ \\
2018 & $50{,}542$ & $8{,}038$ & $15.9\%$ \\
2021 & $35{,}978$ & $6{,}812$ & $18.9\%$ \\
2022 & $22{,}800$ & $5{,}818$ & $25.5\%$ \\
2023 & $17{,}947$ & $7{,}812$ & $43.5\%$ \\
2024 & $10{,}562$ & $5{,}534$ & $52.4\%$ \\
2025 & $2{,}531$ & $1{,}492$ & $59.0\%$ \\
\bottomrule
\end{tabular}
\end{table}

The count series tells a different story. The observed number of AI firms peaks
in 2018 at $8{,}038$ and falls after that. We do not read this as a decline in
AI entrepreneurship. The total observed denominator falls from $54{,}150$ firms
in the 2015 cohort to $2{,}531$ in 2025, because recent firms have not yet had
time to enter commercial registers. The useful fact is therefore not that
observed AI counts fall after 2018, but that AI's share continues to rise even
as the recent denominator becomes incomplete.

\subsection{The trend is not driven by one source}

The rise appears in both commercial databases. Database A rises from $2.7\%$
AI in the 2000 cohort to $63.7\%$ in 2025. Database B rises from $3.6\%$ to
$58.3\%$ over the same window. The levels differ somewhat, but the direction
does not. This matters because a trend visible in only one database could
reflect a vendor-specific tagging or coverage decision. Here, both sources
show the same broad movement.

\subsection{The trend survives a stricter selection rule}

Recent cohorts are not only smaller; they are also more selected. A young firm
is especially likely to appear in a commercial database if it has already done
something visible, such as raising financing. If AI firms are more likely to be
visible early, the rising AI share could partly reflect selection into the
database rather than startup formation itself.

We check this by holding the selection rule fixed: among firms with recorded
financing, the AI share still rises from $5.9\%$ in 2005 to $15.6\%$ in 2015,
$22.0\%$ in 2021, and $68.3\%$ in 2025. The level is higher because financed
firms are more AI-intensive, but the time pattern remains. This does not solve
all selection problems, but it rules out the simplest concern that the trend is
only an artifact of looser recent coverage.

\subsection{Commercial databases miss AI-heavy firms}

The formation series above uses the commercial population because it has the
founding-year metadata needed for a time trend. The broader dataset also
reveals what those commercial sources miss. Firms found outside both
commercial databases are $23.8\%$ AI-related, compared with $12.4\%$ in
database A and $10.7\%$ in database B. The unlisted layer is therefore roughly
twice as AI-intensive as the conventional evidence base.

We do not use the unlisted layer to extend the time series, because its
founding-year coverage is thin. Its value is different: it shows the direction
of coverage bias. If the firms missing from commercial databases are more
likely to be AI-related, then commercial-only studies will understate AI entry.
Section~\ref{sec:checks} tests whether that gap survives alternative AI
definitions, source exclusions, independent labels, and classifier-error
correction.

\subsection{The global pattern is broad growth with cautious geography}

AI intensity rises across every major region we can measure. Between the
2010--2016 and 2020--2025 cohorts, North America's AI share rises from $9.1\%$
to $29.2\%$, Europe's from $7.3\%$ to $21.4\%$, East Asia's from $12.3\%$ to
$25.1\%$, and South and Southeast Asia's from $7.2\%$ to $18.2\%$. Among large
country samples, recent AI intensity is highest in Israel ($37.2\%$), the
United States ($29.4\%$), China ($28.5\%$), Canada ($27.3\%$), and South Korea
($26.1\%$).

The concentration series is more delicate. Through 2017, AI entrepreneurship
appears to become more geographically dispersed: the number of observed
countries with AI firms rises from $59$ in 2000 to $103$ in 2017, and the U.S.\
share of observed AI firms falls from $41.3\%$ to $36.8\%$. After 2020 the
series reverses sharply, with the U.S.\ reaching $85.0\%$ of observed AI firms
in the 2025 cohort. We treat that reversal as a coverage warning rather than a
substantive geography result, because the newest non-U.S.\ firms are indexed
more slowly.

\subsection{AI growth is broad across sectors}

The sector pattern is also one of broad intensification rather than a single
vertical driving the result. Data and analytics rises from $35.9\%$ AI in the
2010--2016 cohort to $71.4\%$ in 2020--2025. Biotech and pharma rises from
$23.2\%$ to $60.0\%$, software and IT from $9.5\%$ to $29.3\%$, and healthcare
from $8.9\%$ to $26.6\%$. Lower-intensity sectors also rise: food and agtech,
consumer and lifestyle, and e-commerce and retail all more than double their
AI shares. The ordering changes less than the levels. Sectors that were already
AI-intensive remain near the top, while nearly every sector moves upward.

% ============================================================================

\section{Findings: Global AI Startup Formation}
\label{sec:findings_global_formation}

The main descriptive pattern is simple: AI has become a much larger share of
new startup formation. The harder part is interpretation. Recent founding
cohorts are still being indexed, so raw counts understate recent formation.
For that reason, the formation result is stated primarily as a \emph{share of
observed firms founded in the same year}, not as the raw number of AI firms.

\subsection{AI's share of new firms rises sharply}

In the commercial population, where founding years are nearly complete, the AI
share rises from $3.0\%$ in 2000 to $10.6\%$ in 2015, $18.9\%$ in 2021, and
$59.0\%$ in 2025 (Table~\ref{tab:formation_selected_years}). The largest break
in the series is between the 2022 and 2023 founding cohorts, when the AI share
jumps from $25.5\%$ to $43.5\%$. That timing is consistent with the public
release of consumer generative AI in late 2022, but we do not interpret the
jump causally. The same cohorts are also affected by indexing lag, and the
paper does not have a design that separates the technology shock from the
change in observability.

\begin{table}[htbp]
\centering\small
\caption{AI share of observed startup formation, selected founding cohorts.
The sample is the commercial population, where founding years are nearly
complete.}
\label{tab:formation_selected_years}
\begin{tabular}{@{}rrrr@{}}
\toprule
\textbf{Founded year} & \textbf{Observed firms} & \textbf{AI firms} &
\textbf{AI share} \\
\midrule
2000 & $28{,}963$ & $874$ & $3.0\%$ \\
2010 & $41{,}513$ & $1{,}996$ & $4.8\%$ \\
2015 & $54{,}150$ & $5{,}764$ & $10.6\%$ \\
2018 & $50{,}542$ & $8{,}038$ & $15.9\%$ \\
2021 & $35{,}978$ & $6{,}812$ & $18.9\%$ \\
2022 & $22{,}800$ & $5{,}818$ & $25.5\%$ \\
2023 & $17{,}947$ & $7{,}812$ & $43.5\%$ \\
2024 & $10{,}562$ & $5{,}534$ & $52.4\%$ \\
2025 & $2{,}531$ & $1{,}492$ & $59.0\%$ \\
\bottomrule
\end{tabular}
\end{table}

The count series tells a different story. The observed number of AI firms peaks
in 2018 at $8{,}038$ and falls after that. We do not read this as a decline in
AI entrepreneurship. The total observed denominator falls from $54{,}150$ firms
in the 2015 cohort to $2{,}531$ in 2025, because recent firms have not yet had
time to enter commercial registers. The useful fact is therefore not that
observed AI counts fall after 2018, but that AI's share continues to rise even
as the recent denominator becomes incomplete.

\subsection{The trend is not driven by one source}

The rise appears in both commercial databases. Database A rises from $2.7\%$
AI in the 2000 cohort to $63.7\%$ in 2025. Database B rises from $3.6\%$ to
$58.3\%$ over the same window. The levels differ somewhat, but the direction
does not. This matters because a trend visible in only one database could
reflect a vendor-specific tagging or coverage decision. Here, both sources
show the same broad movement.

\subsection{The trend survives a stricter selection rule}

Recent cohorts are not only smaller; they are also more selected. A young firm
is especially likely to appear in a commercial database if it has already done
something visible, such as raising financing. If AI firms are more likely to be
visible early, the rising AI share could partly reflect selection into the
database rather than startup formation itself.

We check this by holding the selection rule fixed: among firms with recorded
financing, the AI share still rises from $5.9\%$ in 2005 to $15.6\%$ in 2015,
$22.0\%$ in 2021, and $68.3\%$ in 2025. The level is higher because financed
firms are more AI-intensive, but the time pattern remains. This does not solve
all selection problems, but it rules out the simplest concern that the trend is
only an artifact of looser recent coverage.

\subsection{Commercial databases miss AI-heavy firms}

The formation series above uses the commercial population because it has the
founding-year metadata needed for a time trend. The broader dataset also
reveals what those commercial sources miss. Firms found outside both
commercial databases are $23.8\%$ AI-related, compared with $12.4\%$ in
database A and $10.7\%$ in database B. The unlisted layer is therefore roughly
twice as AI-intensive as the conventional evidence base.

We do not use the unlisted layer to extend the time series, because its
founding-year coverage is thin. Its value is different: it shows the direction
of coverage bias. If the firms missing from commercial databases are more
likely to be AI-related, then commercial-only studies will understate AI entry.
Section~\ref{sec:checks} tests whether that gap survives alternative AI
definitions, source exclusions, independent labels, and classifier-error
correction.

\subsection{The global pattern is broad growth with cautious geography}

AI intensity rises across every major region we can measure. Between the
2010--2016 and 2020--2025 cohorts, North America's AI share rises from $9.1\%$
to $29.2\%$, Europe's from $7.3\%$ to $21.4\%$, East Asia's from $12.3\%$ to
$25.1\%$, and South and Southeast Asia's from $7.2\%$ to $18.2\%$. Among large
country samples, recent AI intensity is highest in Israel ($37.2\%$), the
United States ($29.4\%$), China ($28.5\%$), Canada ($27.3\%$), and South Korea
($26.1\%$).

The concentration series is more delicate. Through 2017, AI entrepreneurship
appears to become more geographically dispersed: the number of observed
countries with AI firms rises from $59$ in 2000 to $103$ in 2017, and the U.S.\
share of observed AI firms falls from $41.3\%$ to $36.8\%$. After 2020 the
series reverses sharply, with the U.S.\ reaching $85.0\%$ of observed AI firms
in the 2025 cohort. We treat that reversal as a coverage warning rather than a
substantive geography result, because the newest non-U.S.\ firms are indexed
more slowly.

\subsection{AI growth is broad across sectors}

The sector pattern is also one of broad intensification rather than a single
vertical driving the result. Data and analytics rises from $35.9\%$ AI in the
2010--2016 cohort to $71.4\%$ in 2020--2025. Biotech and pharma rises from
$23.2\%$ to $60.0\%$, software and IT from $9.5\%$ to $29.3\%$, and healthcare
from $8.9\%$ to $26.6\%$. Lower-intensity sectors also rise: food and agtech,
consumer and lifestyle, and e-commerce and retail all more than double their
AI shares. The ordering changes less than the levels. Sectors that were already
AI-intensive remain near the top, while nearly every sector moves upward.
